\subsection{Aliasing Error}
\label{SEC:Aliasing}

%%%%%%%%%%%%%%%%%%%%%%%%%%%%%%%%%%%%%%%%%%%%%%%%%%%%%%%%%%%%%%%%%%%%%%%%%%%%%%%%

\key*{Aliasing error}[aliasing error] occurs when the polynomial basis
functions are insufficient to represent the nonlinear terms.
%
Consider the simplified example shown in \fref{FIG:Aliasing},
in which the exact solution $u$ is given by a fourth order polynomial
and the exact flux $u^2$ is given by an eighth order polynomial.
%
In this case, a fourth order polynomial for the numerical solution
can be used to exactly represent the physical solution.
%
Evaluating the flux at the solution points will also result in a
fourth order polynomial that matches the exact flux at those points.
%
However, at any other location, the fourth order polynomial flux
is unable to accurately approximate the eighth order exact flux.
%
The area between the exact flux and polynomial flux curves represents
the aliasing error.

%%%%%%%%%%%%%%%%%%%%%%%%%%%%%%%%%%%%%%%%%%%%%%%%%%%%%%%%%%%%%%%%%%%%%%%%%%%%%%%%

\begin{figure}[htb]
  \begin{centering}
    \insertimage[width=4.0in]{Aliasing/Aliasing}
    \caption{Example of aliasing error.}
    \label{FIG:Aliasing}
  \end{centering}
\end{figure}

%%%%%%%%%%%%%%%%%%%%%%%%%%%%%%%%%%%%%%%%%%%%%%%%%%%%%%%%%%%%%%%%%%%%%%%%%%%%%%%%

\Linecite{Spiegel_2015b} %AIAA-2015-2744
explored one means of reducing aliasing error
using an early version of GFR.
%
Instead of directly computing the flux using the same polynomial
order and Gauss points as the solution, the flux is evaluated
using a higher quadrature order
(\thatis a higher order polynomial and its corresponding Gauss points).
%
This technique of over-integration by over-sampling at the quarature points
requires that the solution polynomial for the conserved variables
be evaluated (\thatis interpolated) at the additional set of quadrature points
in order to use those values in computing the flux.
%
Then that higher order quadrature polynomial is projected
(in a least squares sense) onto a flux polynomial of the same order
as the solution.
%
This reduces the {\em overall error} between the exact flux and
polynomial flux curves, but in doing so the {\em specific error}
between the exact flux and polynomial flux curves evaluated at the
solution Gauss points will now be non-zero.

%%%%%%%%%%%%%%%%%%%%%%%%%%%%%%%%%%%%%%%%%%%%%%%%%%%%%%%%%%%%%%%%%%%%%%%%%%%%%%%%

The main disadvantage to reducing aliasing error in this manner is
the additional cost involved in supersampling at the additional
\key{quadrature points}[points!quadrature] and then projecting that
polynomial down to the solution order.
%
For that reason, subsequent development of GFR (\thatis the viscous terms)
has not incorporated the quadrature evaluation.

%%%%%%%%%%%%%%%%%%%%%%%%%%%%%%%%%%%%%%%%%%%%%%%%%%%%%%%%%%%%%%%%%%%%%%%%%%%%%%%%

